\documentclass[12pt, a4paper]{article}

\usepackage[margin=1.8cm,top=2cm,bottom=2cm]{geometry}
\usepackage{amsmath,amssymb,amsfonts}
\usepackage{graphicx}
\usepackage[colorlinks=true,linkcolor=blue,citecolor=blue,urlcolor=blue]{hyperref}
\usepackage{bm,xcolor}
\usepackage[numbers,sort&compress]{natbib}
\usepackage{float}

\newcommand{\Geff}{G_{\mathrm{eff}}}
\newcommand{\cs}{c_s}
\newcommand{\meff}{m_{\mathrm{eff}}}
\newcommand{\heff}{\hbar_{\mathrm{eff}}}
\newcommand{\lP}{\ell_{\mathrm{P}}}
\newcommand{\mP}{m_{\mathrm{P}}}

\begin{document}

\title{Emergent Electromagnetism from the Gauged\\
Logarithmic Superfluid Vacuum}

\author{B. B. Kulangiev\\\small{Haifa, Israel}}
\date{April 2026}
\maketitle

\begin{abstract}
\noindent
The relativistic logarithmic Klein-Gordon equation
for the vacuum condensate possesses a global $U(1)$
phase symmetry. We show that promoting this to a
local gauge symmetry---the minimal coupling
$\partial_\mu S \to \partial_\mu S - eA_\mu$---produces
Maxwell's equations as the Euler-Lagrange equations
for the gauge field $A_\mu$, with the condensate
current $J^\nu = (e\meff/\heff^2)\,\rho\,u^\nu$
as source. The Madelung decomposition of the gauged
equation yields the acoustic metric equations
(emergent gravity) and Maxwell's equations (emergent
electromagnetism) simultaneously from a single
Lagrangian. Gravity is sourced by the soliton energy
density (mass); electromagnetism is sourced by the
topological winding number of vortex-Gausson defects,
which quantizes electric charge as $Q = ne$ with
$n \in \mathbb{Z}$. Both fields propagate at the
sound speed $c$ on the same acoustic metric,
automatically ensuring $c_{\mathrm{grav}} =
c_{\mathrm{EM}} = c$ as observed. The London
penetration depth of the condensate determines the
electromagnetic screening length; in the electrically
neutral vacuum background, it diverges, permitting
free propagation of electromagnetic waves. Near
vortex cores, the screening becomes significant,
producing the Coulomb field as a condensate boundary
effect. The framework unifies gravity and
electromagnetism as two aspects of the same
condensate dynamics---the acoustic metric and the
gauge connection---while preserving their distinct
physical origins: mass versus topology.
\end{abstract}


% ====================================================================
\section{Introduction}
\label{sec:intro}
% ====================================================================

In previous work, we showed that the Madelung
decomposition of the relativistic logarithmic
Klein-Gordon equation produces an acoustic metric
equivalent to the Schwarzschild geometry in the
weak-field limit, yielding the PPN parameter
$\gamma = 1$ and an emergent gravitational constant
$\Geff = \eta c^2/(4\pi\rho_0)$. The vacuum
condensate's density $\rho$ and flow velocity
$u_\mu = \partial_\mu S/\meff$ generate the
gravitational field; topological defects
(vortex-Gaussons) play the role of massive particles.

This framework has a conspicuous unused structure.
The logarithmic wave equation
\begin{equation}
    \Box\psi + \mu^2\psi
    - b\,\psi\,\ln\!\left(\frac{|\psi|^2}{\rho_0}\right)
    = 0
\label{eq:log_kg}
\end{equation}
is invariant under the global phase rotation
$\psi \to \psi\,e^{i\alpha}$ for constant $\alpha$.
This $U(1)$ symmetry is associated, via Noether's
theorem, with the conservation of the condensate
particle number (the continuity
equation $\partial_\mu(\rho\,\partial^\mu S) = 0$).

In the Standard Model, the fundamental interactions
arise from \emph{gauging} internal symmetries:
promoting global invariances to local ones and
introducing gauge fields to maintain covariance.
Electromagnetism itself originates from gauging the
$U(1)$ phase symmetry of the electron
field~\cite{weinberg_1995}.

In this paper, we apply the same principle to the
vacuum condensate. Promoting the global $U(1)$ of
Eq.~(\ref{eq:log_kg}) to a local symmetry requires
the introduction of a gauge field $A_\mu$ and the
minimal coupling $\partial_\mu \to D_\mu =
\partial_\mu - i(e/\hbar)A_\mu$. We show that the
Euler-Lagrange equations for $A_\mu$ are precisely
Maxwell's equations, with the condensate current as
source. Together with the acoustic metric (gravity),
this produces a unified description of two
fundamental interactions from a single condensate
Lagrangian.


% ====================================================================
\section{The Gauged Lagrangian}
\label{sec:lagrangian}
% ====================================================================

The ungauged Lagrangian density for the logarithmic
condensate is:
\begin{equation}
    \mathcal{L}_0 = (\partial_\mu\psi^*)(\partial^\mu\psi)
    - \mu^2|\psi|^2
    + b\,|\psi|^2\ln\!\left(\frac{|\psi|^2}{\rho_0}\right)\,.
\label{eq:L0}
\end{equation}
This is invariant under $\psi \to \psi\,e^{i\alpha}$
for constant $\alpha$.

Promoting to local gauge invariance
$\psi \to \psi\,e^{i\alpha(x)}$ requires replacing
$\partial_\mu \to D_\mu = \partial_\mu
- i(e/\hbar)A_\mu$ and adding the gauge field
kinetic term. The gauged Lagrangian is:
\begin{equation}
    \boxed{\;\mathcal{L} = (D_\mu\psi)^*(D^\mu\psi)
    - \mu^2|\psi|^2
    + b\,|\psi|^2\ln\!\left(\frac{|\psi|^2}{\rho_0}\right)
    - \frac{1}{4}F_{\mu\nu}F^{\mu\nu}\;}
\label{eq:L_gauged}
\end{equation}
where $F_{\mu\nu} = \partial_\mu A_\nu
- \partial_\nu A_\mu$ is the field strength tensor
and $e$ is the gauge coupling constant.

The term $-\frac{1}{4}F_{\mu\nu}F^{\mu\nu}$ is the
unique Lorentz-invariant, gauge-invariant,
parity-preserving kinetic term for $A_\mu$ with at
most two derivatives~\cite{weinberg_1995}. Its
inclusion is not optional: without it, $A_\mu$ would
have no dynamics and could be gauged away.


% ====================================================================
\section{Madelung Decomposition of the Gauged Equation}
\label{sec:madelung}
% ====================================================================

\subsection{The gauge-invariant velocity}

Applying the Madelung transformation
$\psi = \sqrt{\rho}\,e^{iS/\heff}$, the covariant
derivative becomes:
\begin{equation}
    D_\mu\psi = \left(\frac{\partial_\mu\sqrt{\rho}}
    {\sqrt{\rho}} + \frac{i}{\heff}
    (\partial_\mu S - eA_\mu)\right)\psi\,.
\end{equation}
The gauge-invariant superfluid four-velocity is:
\begin{equation}
    \boxed{\;u_\mu = \frac{1}{\meff}
    (\partial_\mu S - eA_\mu)\;}
\label{eq:u_gauge}
\end{equation}
This is the central modification. In the ungauged
theory, $u_\mu = \partial_\mu S/\meff$; the gauge
field shifts the velocity by $-eA_\mu/\meff$. This
is the relativistic generalization of the London
relation in superconductors: $m\mathbf{v}_s =
\hbar\nabla\phi - e\mathbf{A}$~\cite{london_1935}.

\newpage
\subsection{The modified fluid equations}

Separating the Euler-Lagrange equation for $\psi$
into real and imaginary parts yields the gauged
Madelung equations:

\emph{Continuity:}
\begin{equation}
    \partial_\mu(\rho\,u^\mu) = 0\,.
\label{eq:continuity_gauged}
\end{equation}

\emph{Hamilton-Jacobi:}
\begin{equation}
    \frac{\meff^2}{\heff^2}\,u_\mu u^\mu
    = \mu^2 - b\,\ln\!\left(\frac{\rho}{\rho_0}\right)
    + Q[\rho]\,,
\label{eq:hj_gauged}
\end{equation}
where $Q = \Box\sqrt{\rho}/\sqrt{\rho}$ is the
quantum potential (Bohm potential), unchanged from
the ungauged case.

These equations are \emph{identical in form} to the
ungauged Madelung equations, with the single
replacement $\partial_\mu S \to \partial_\mu S
- eA_\mu$. The acoustic metric depends on $\rho$
and $u_\mu$, both of which are gauge-invariant. The
entire gravitational sector---the PPN parameter
$\gamma = 1$, the emergent $\Geff$, the
Painlev\'{e}-Gullstrand flow---is preserved without
modification.


% ====================================================================
\section{Maxwell's Equations from the Gauge Field Variation}
\label{sec:maxwell}
% ====================================================================

The Euler-Lagrange equation for $A_\mu$ is obtained
by varying the Lagrangian~(\ref{eq:L_gauged}) with
respect to $A^\nu$.

The condensate contribution:
\begin{equation}
    \frac{\delta}{\delta A^\nu}
    \left[(D_\mu\psi)^*(D^\mu\psi)\right]
    = -\frac{e\,\meff}{\heff^2}\,\rho\,u^\nu\,.
\end{equation}
The gauge field contribution:
\begin{equation}
    \frac{\delta}{\delta A^\nu}
    \left[-\frac{1}{4}F_{\alpha\beta}F^{\alpha\beta}\right]
    = \partial_\mu F^{\mu\nu}\,.
\end{equation}
Setting the total variation to zero:
\begin{equation}
    \boxed{\;\partial_\mu F^{\mu\nu} = J^\nu\;}
\label{eq:maxwell}
\end{equation}
where the four-current is:
\begin{equation}
    J^\nu = \frac{e\,\meff}{\heff^2}\,\rho\,u^\nu\,.
\label{eq:current}
\end{equation}
Equation~(\ref{eq:maxwell}) is Maxwell's inhomogeneous
equation. The homogeneous equations
$\partial_{[\mu}F_{\nu\lambda]} = 0$ follow
identically from the definition
$F_{\mu\nu} = \partial_\mu A_\nu
- \partial_\nu A_\mu$ (the Bianchi identity).

In three-vector notation:
\begin{align}
    \nabla\cdot\mathbf{E} &= \rho_{\mathrm{ch}}/\epsilon_0
    && \text{(Gauss)}\,,\label{eq:gauss}\\
    \nabla\cdot\mathbf{B} &= 0
    && \text{(no monopoles)}\,,\\
    \nabla\times\mathbf{E} &= -\frac{\partial\mathbf{B}}
    {\partial t}
    && \text{(Faraday)}\,,\\
    \nabla\times\mathbf{B} &= \mu_0\mathbf{J}
    + \mu_0\epsilon_0\frac{\partial\mathbf{E}}{\partial t}
    && \text{(Amp\`{e}re-Maxwell)}\,,\label{eq:ampere}
\end{align}
where $\rho_{\mathrm{ch}} = J^0$ is the charge
density and $\mathbf{J} = (J^1, J^2, J^3)$ is the
current density. These are the full Maxwell equations,
derived from the gauge field variation of the
condensate Lagrangian.


% ====================================================================
\section{Charge Quantization from Vortex Topology}
\label{sec:charge}
% ====================================================================

The vortex-Gausson solitons of the logarithmic
condensate have quantized phase winding. For a vortex
of winding number $n$, the phase increases by
$2\pi n\heff$ around any closed loop encircling the
core:
\begin{equation}
    \oint \partial_\mu S\,dx^\mu = 2\pi n\heff\,,
    \qquad n \in \mathbb{Z}\,.
\label{eq:winding}
\end{equation}
The gauge-invariant circulation is:
\begin{equation}
    \oint u_\mu\,dx^\mu
    = \frac{1}{\meff}\oint(\partial_\mu S
    - eA_\mu)\,dx^\mu
    = \frac{2\pi n\heff}{\meff}
    - \frac{e}{\meff}\oint A_\mu\,dx^\mu\,.
\label{eq:circulation}
\end{equation}
In the absence of an external gauge field
($A_\mu = 0$ at infinity), the total charge enclosed
by the loop is determined by the winding number:
\begin{equation}
    \boxed{\;Q = n\,e\,, \qquad n \in \mathbb{Z}\;}
\label{eq:charge_quantized}
\end{equation}
Electric charge is quantized because vortex winding
numbers are integers. This is the same mechanism by
which magnetic flux is quantized in
superconductors~\cite{london_1950}: the single-valuedness
of the order parameter forces the enclosed quantity
to be an integer multiple of the fundamental quantum.

The physical identification is:
\begin{itemize}
\item Vortex-Gaussons with $n = +1$: positively
  charged particles (charge $+e$).
\item Anti-vortices with $n = -1$: negatively charged
  particles (charge $-e$).
\item Pure Gaussons with $n = 0$: electrically neutral
  particles (charge $0$).
\end{itemize}
The vortex-Gausson thus carries both mass
(from its soliton energy $E_v$, generating gravity)
and charge (from its winding number $n$, generating
electromagnetism).


% ====================================================================
\section{Electromagnetic Waves in the Condensate}
\label{sec:em_waves}
% ====================================================================

In the vacuum background ($\rho = \rho_0$,
$u_\mu = 0$, no net charge), the linearized Maxwell
equations~(\ref{eq:maxwell}) reduce to:
\begin{equation}
    \Box A_\mu = 0
    \qquad\text{(in Lorenz gauge)}\,,
\label{eq:em_wave}
\end{equation}
where $\Box = \partial_t^2/c^2 - \nabla^2$ and
$c = c_s(\rho_0)$ is the sound speed at the
background density. Electromagnetic waves propagate
at the speed of sound in the condensate.

This provides a unified explanation: the speed of
light is the speed of sound because the photon is
a gauge excitation propagating on the same condensate
whose acoustic properties define $c$. The identity
$c_{\mathrm{EM}} = c_{\mathrm{grav}} = c_{\mathrm{sound}}$
is not a coincidence but a structural consequence
of both fields emerging from the same substrate.

Near a gravitating mass, the acoustic metric
modifies the effective spacetime seen by all
excitations equally. Electromagnetic waves are
deflected, redshifted, and time-delayed by the
Painlev\'{e}-Gullstrand flow, reproducing the
standard general relativistic predictions. The
Shapiro time delay, the gravitational redshift, and
the bending of light are automatically consistent
with $\gamma = 1$ because the photon (gauge
excitation) and the phonon (density excitation)
share the same acoustic metric.


% ====================================================================
\section{The London Equation and Electromagnetic Screening}
\label{sec:london}
% ====================================================================

In a region where the condensate density is uniform
($\rho = \rho_0$) and the phase is constant
($\partial_\mu S = 0$), the gauge-invariant
velocity~(\ref{eq:u_gauge}) reduces to:
\begin{equation}
    u_\mu = -\frac{e}{\meff}\,A_\mu\,.
\label{eq:london_velocity}
\end{equation}
Substituting into the current~(\ref{eq:current}):
\begin{equation}
    J^\nu = -\frac{e^2\rho_0}{\meff\heff^2}\,A^\nu\,.
\label{eq:london_current}
\end{equation}
Combined with Maxwell's equation~(\ref{eq:maxwell}),
this gives the Proca-like equation:
\begin{equation}
    \Box A^\nu + \frac{1}{\lambda_L^2}\,A^\nu = 0\,,
\label{eq:proca}
\end{equation}
where the London penetration depth is:
\begin{equation}
    \lambda_L = \frac{\heff}{e}\,
    \frac{1}{\sqrt{\rho_0}}\,.
\label{eq:london_depth}
\end{equation}
Inside a region of uniform condensate with net
charge density, electromagnetic fields are
exponentially screened on the scale $\lambda_L$.
This is the Meissner effect: the condensate expels
electromagnetic fields from its bulk.

In the physical vacuum, however, the background
condensate is electrically neutral. There is no net
winding ($\langle n \rangle = 0$), and the
macroscopic phase $S$ is not locked to a constant---it
varies spatially (producing the PG flow for
gravity). The screening
condition~(\ref{eq:london_velocity}) does not apply
globally, and electromagnetic waves propagate
freely, as observed.

The London screening becomes significant only in
the immediate vicinity of a vortex-Gausson core
($r \lesssim \xi$), where the condensate density
drops to zero and the phase winds rapidly. The
transition from the screened interior to the
unscreened bulk produces the Coulomb field: the
$1/r$ potential is a boundary effect between the
vortex core (where $A_\mu$ is confined) and the
free vacuum (where $A_\mu$ propagates).


% ====================================================================
\section{Unified Picture: Gravity and Electromagnetism}
\label{sec:unified}
% ====================================================================

The gauged logarithmic condensate produces two
emergent interactions from a single Lagrangian:

\emph{Gravity} arises from the acoustic metric
$g_{\mu\nu}(\rho, u)$. Its source is the soliton
energy density (mass). Its field equation is the
Bernoulli relation and its consequence, the
Painlev\'{e}-Gullstrand acoustic metric equivalent
to Schwarzschild. Its coupling constant is
$G = \eta c^2/(4\pi\rho_0)$.

\emph{Electromagnetism} arises from the gauge field
$A_\mu$ required by local $U(1)$ invariance. Its
source is the topological winding number (charge).
Its field equation is Maxwell's
$\partial_\mu F^{\mu\nu} = J^\nu$. Its coupling
constant is $e$.

The two interactions are coupled through the
condensate:
\begin{itemize}
\item Charged particles (vortex-Gaussons) create
  both gravitational (acoustic metric) and
  electromagnetic (gauge) fields simultaneously.
\item The electromagnetic field modifies the
  condensate flow through the London
  relation~(\ref{eq:london_velocity}).
\item The gravitational flow (PG inflow) modifies
  the effective spacetime for electromagnetic
  wave propagation---this is the gravitational
  deflection of light.
\end{itemize}

Despite this coupling, the two interactions have
fundamentally different origins within the
condensate. Gravity is universal (all solitons
have energy, therefore mass), while
electromagnetism is topological (only vortex
solitons with $n \neq 0$ carry charge). This
explains why gravity couples to all matter while
electromagnetism couples only to charged matter---a
distinction that is postulated in GR but derived
here from the condensate structure.

% ====================================================================
\section{Discussion}
\label{sec:discussion}
% ====================================================================

\subsection{What is derived and what is assumed}

The derivation rests on three inputs:

(i) The logarithmic wave equation~(\ref{eq:log_kg}),
which defines the condensate dynamics.

(ii) The requirement of local $U(1)$ gauge
invariance, which forces the introduction of $A_\mu$.

(iii) The gauge field kinetic term
$-\frac{1}{4}F_{\mu\nu}F^{\mu\nu}$, which is the
unique allowed term.

Of these, (i) is the foundation of the entire
framework. Input (ii) is a symmetry principle---the
same one used in the Standard Model to derive
electromagnetism from the electron field's $U(1)$
symmetry~\cite{weinberg_1995}. Input (iii) is
uniquely determined by Lorentz invariance, gauge
invariance, and the requirement of at most two
derivatives.

The coupling constant $e$ is \emph{not} determined
by the condensate parameters. Deriving its value
(equivalently, the fine structure constant
$\alpha = e^2/(4\pi\epsilon_0\hbar c) \approx
1/137$) would require computing the self-energy
of a vortex-Gausson in the gauged condensate and
matching it to the observed electron mass. This
remains an open problem.

\subsection{Comparison with the Standard Model}

In the Standard Model, electromagnetism arises from
gauging the $U(1)_Y$ hypercharge symmetry of the
fermion and Higgs fields. The photon is the massless
gauge boson associated with the unbroken $U(1)_{\mathrm{EM}}$
after electroweak symmetry breaking. The electron's
charge is a coupling constant.

In the present framework, electromagnetism arises
from gauging the $U(1)$ phase symmetry of the vacuum
condensate. The photon is the massless gauge
excitation propagating on the condensate. The
electron's charge is the same coupling constant $e$,
and charge quantization follows from vortex topology
rather than from anomaly cancellation.

The key structural difference is that here, the
condensate provides both the gravitational field
(acoustic metric) and the electromagnetic field
(gauge connection), whereas the Standard Model
treats gravity as external.

\subsection{Limitations}

The framework produces only a $U(1)$ gauge field.
The Standard Model requires $SU(3) \times SU(2)
\times U(1)$. Extending to non-Abelian gauge groups
would require a multi-component order parameter
(e.g., a spinor or matrix-valued condensate), which
is beyond the scope of the present single-component
model.

The identification of electric charge with vortex
winding number predicts that all charged particles
have integer multiples of $e$. This is consistent
with observation (quarks carry fractional charges,
but they are confined, and all observable states
carry integer charges). However, deriving quark
confinement from the condensate would require the
$SU(3)$ extension mentioned above.

\subsection{Speculative application: planetary magnetic fields}

The axial outflow derived in the gravitational
sector~\cite{kulangiev_2026b} has a potential
electromagnetic consequence within the present
framework. The outflow along the rotation axis
creates a preferred direction that may align the
vortex axes of embedded solitons. A population of
$N$ vortex-Gaussons with correlated winding axes
constitutes a net current loop. By Amp\`{e}re's
law~(\ref{eq:ampere}), a current along an axis
produces a dipole magnetic field. If this mechanism
operates, the predicted scaling is:
\begin{equation}
    |\mathbf{B}_{\mathrm{dipole}}| \propto
    \Omega\,M\,f_{\mathrm{liquid}}\,,
\label{eq:B_scaling}
\end{equation}
where $\Omega$ is the rotation rate (setting the
outflow strength), $M$ the body mass (setting the PG
inflow strength), and $f_{\mathrm{liquid}}$ the
fraction of the core that supports free vortex motion
(liquid vs.\ solid).

This qualitatively explains the observed pattern of
planetary magnetic fields:
\begin{center}
\begin{tabular}{lcccl}
\hline
Body & $\Omega$ & Core & $B$ & Prediction \\
\hline
Jupiter & fast & metallic H & strongest & \checkmark \\
Earth & moderate & liquid Fe & moderate & \checkmark \\
Mercury & slow & partial liq. & weak & \checkmark \\
Venus & very slow & liquid & $\approx 0$ & \checkmark \\
Mars & moderate & solid & $\approx 0$ & \checkmark \\
\hline
\end{tabular}
\end{center}
Venus fails despite having a liquid core because its
rotation rate is too slow to generate significant
axial outflow. Mars fails despite moderate rotation
because its core has solidified, preventing vortex
alignment. A quantitative treatment requires coupling
the vacuum flow equations to magnetohydrodynamic core
models and is deferred to future work.

% ====================================================================
\section{Conclusion}
\label{sec:conclusion}
% ====================================================================

We have shown that gauging the $U(1)$ phase symmetry
of the logarithmic superfluid vacuum produces
Maxwell's equations alongside the acoustic metric
equations from a single Lagrangian. The Madelung
decomposition of the gauged equation separates
cleanly into:

(1) The gravitational sector: the acoustic metric
$g_{\mu\nu}(\rho, u)$, governed by the Bernoulli
relation and the Painlev\'{e}-Gullstrand flow,
sourced by mass (soliton energy).

(2) The electromagnetic sector: the gauge field
$A_\mu$, governed by Maxwell's equations
$\partial_\mu F^{\mu\nu} = J^\nu$, sourced by
charge (vortex winding number).

Electric charge is quantized ($Q = ne$, $n \in
\mathbb{Z}$) because vortex winding numbers are
integers. Both fields propagate at the sound speed
$c$, explaining the identity
$c_{\mathrm{EM}} = c_{\mathrm{grav}}$ as a
structural consequence rather than a coincidence.
The framework predicts that axial outflow of the
Painlev\'{e}-Gullstrand flow aligns embedded
vortex-Gaussons, producing dipole magnetic fields
correlated with rotation rate and core
state---qualitatively consistent with the observed
pattern of planetary magnetism.

The coupling constant $e$ (the fine structure
constant) is not determined by the condensate
parameters and remains an open problem. Extension
to non-Abelian gauge groups requires a
multi-component condensate and is deferred to
future work.


% ====================================================================
\newpage
\section*{Acknowledgments}
% ====================================================================

The author thanks K.~G.~Zloshchastiev for foundational
work on the logarithmic superfluid vacuum.
Claude (Anthropic, claude-opus-4-6) and Gemini (Google, Gemini 3.1 Pro) were used as editorial and computational tools during manuscript preparation. All scientific content, equations, and
interpretations are the author's own.

\begin{thebibliography}{20}

\bibitem{kulangiev_2026b}
B.~B.~Kulangiev,
Emergent Newtonian Gravity and the Gravitational
Feynman-Onsager Relation in the Logarithmic Superfluid Vacuum,
Preprint (2026).
\url{https://kulangiev.github.io#combined}

\bibitem{zloshchastiev_2023}
K.~G.~Zloshchastiev,
Derivation of emergent spacetime metric, gravitational
potential and speed of light in superfluid vacuum theory,
Universe \textbf{9}, 234 (2023).

\bibitem{zloshchastiev_2020}
K.~G.~Zloshchastiev,
An alternative to dark matter and dark energy:
Scale-dependent gravity in superfluid vacuum theory,
Universe \textbf{6}, 180 (2020).

\bibitem{zloshchastiev_2011}
K.~G.~Zloshchastiev,
Spontaneous symmetry breaking and mass generation as
built-in phenomena in logarithmic nonlinear quantum
theory,
Acta Phys.\ Polon.\ B \textbf{42}, 261 (2011).

\bibitem{bialynicki_birula_1979}
I.~Bia{\l}ynicki-Birula and J.~Mycielski,
Gaussons: Solitons of the logarithmic Schr\"{o}dinger
equation,
Phys.\ Scr.\ \textbf{20}, 539 (1979).

\bibitem{weinberg_1995}
S.~Weinberg,
\emph{The Quantum Theory of Fields}, Vol.~2
(Cambridge, 1996).

\bibitem{london_1935}
F.~London and H.~London,
The electromagnetic equations of the supraconductor,
Proc.\ Roy.\ Soc.\ A \textbf{149}, 71 (1935).

\bibitem{london_1950}
F.~London,
\emph{Superfluids}, Vol.~1 (Wiley, 1950).

\bibitem{volovik_2003}
G.~E.~Volovik,
\emph{The Universe in a Helium Droplet} (Oxford, 2003).

\bibitem{visser_1998}
M.~Visser,
Acoustic black holes: Horizons, ergospheres and Hawking
radiation,
Class.\ Quant.\ Grav.\ \textbf{15}, 1767 (1998).

\bibitem{unruh_1981}
W.~G.~Unruh,
Experimental black-hole evaporation?,
Phys.\ Rev.\ Lett.\ \textbf{46}, 1351 (1981).

\end{thebibliography}

\end{document}
